\section{Benchmark-Analyse: Kühlsysteme der FS-Teams 2025}

Im Rahmen der Konzeptentwicklung wurde eine Benchmark-Analyse der Engineering Design Reports (EDR) führender Teams durchgeführt. Dabei zeigten sich zwei grundlegend verschiedene Strategien: die konventionelle Wasserkühlung für hohe Leistungsdichte und die innovative Luftkühlung zur Gewichtsoptimierung.

\begin{table}[h]
    \centering
    \caption{Vergleich der Kühlkonzepte ausgewählter Teams (Saison 2025)}
    \label{tab:kuehlung_benchmark}
    \begin{tabular}{l|p{4cm}|p{4cm}|p{3.5cm}}
        \hline
        \textbf{Team} & \textbf{Kühlung Antrieb} & \textbf{Kühlung Akkumulator} & \textbf{Besonderheiten} \\ \hline
        \textbf{Schanzer Racing} & \textbf{Wassergekühlt}: 2 Kreisläufe, 4 Radiatoren in den Sidepods & \textbf{Luftgekühlt}: Optimierte Sidepod-Einlässe & Kühlung von 4 Motoren + 2 Dual-Invertern \\ \hline
        \textbf{Joanneum Racing} & \textbf{Luftgekühlt}: Aktive Lüfter ziehen Luft durch Uprights & \textit{Nicht explizit spezifiziert} & Verzicht auf Pumpen/Radiatoren spart Gewicht \\ \hline
        \textbf{Ecurie Aix}      & Fokus auf Verlustminimierung im Inverter/Motor & \textbf{Optimiertes Thermomanagement} & Fokus auf Redesign zur Effizienzsteigerung \\ \hline
        \textbf{AMZ Racing}      & Hydraulische Elemente in Sidepods integriert & \textit{Keine Detailangabe} & Fokus auf Packaging in den Sidepods \\ \hline
        \textbf{Strohm + Söhne}  & AMK FSE Kit Standardkühlung & Fokus auf Rekuperations-Management & Effizienz durch Energiemanagement \\ \hline
    \end{tabular}
\end{table}

\subsection{Detaillierte Betrachtung der Konzepte}

\subsubsection{Aktive Luftkühlung (Beispiel: Joanneum Racing Graz)}
Ein interessanter Ansatz findet sich beim JR25 von Joanneum Racing. Das Team setzt auf eine reine Luftkühlung der Radnabenmotoren. Hierbei sind Kühlrippen direkt in die Uprights integriert, und elektrische Lüfter sorgen für einen konstanten Luftstrom. Der primäre Vorteil liegt in der Reduktion der Systemkomplexität: Da keine Wasserpumpen, Schläuche oder Ausgleichsbehälter benötigt werden, sinkt das Fahrzeuggewicht und das Risiko für Leckagen während des Endurance-Wettbewerbs wird eliminiert.

\subsubsection{Zwei-Kreis-Wasserkühlung (Beispiel: Schanzer Racing Electric)}
Das Team Schanzer Racing (SRe25) verfolgt hingegen ein klassisches Hochleistungskonzept. Um die thermische Last der vier Motoren und der zwei Dual-Inverter (platziert im Front- und Heckbereich) abzuführen, werden zwei separate Wasserkreisläufe verwendet. Die Wärmeabfuhr erfolgt über vier Radiatoren, die in den aerodynamisch optimierten Sidepods untergebracht sind. Dieses System ermöglicht eine stabilere Temperaturführung bei hohen Umgebungstemperaturen, erkauft dies jedoch durch ein höheres Systemgewicht.

\subsubsection{Thermomanagement im Akkumulator (Beispiel: Ecurie Aix)}
Ecurie Aix legt beim eax04 den Schwerpunkt auf die thermische Optimierung des Akkumulators. Durch ein Redesign der Zellkontaktierung und des internen Aufbaus sollen ohmsche Verluste minimiert werden, was die thermische Last von vornherein reduziert. Ziel ist es, die Betriebstemperatur der Zellen auch ohne massiven Kühlaufwand im optimalen Fenster zu halten.